\section{The longest forced wins}\label{sec:longest}

How long can a won game last when both sides play perfectly, the winner as fast as possible and the loser as
slowly as possible? The tables answer this directly, but one more step is needed to get an answer about
real games.

\subsection{Over all positions: 215 plies, in positions no game reaches}

The largest distance in the tables is $\max_x d(x) = 215$. It occurs in four positions (two mirror pairs),
all with the goats to move after 4 or 5 goats are placed and none captured, for example tigers on
$\{4, 12, 13\}$ and goats on $\{3, 5, 6, 22\}$ (Figure~\ref{fig:records}, right). The goats win, and they have
exactly one winning move.

None of the four can occur in a game. This was noticed by looking at the board, and the argument is short.

\begin{lemma}\label{lem:unreachable}
Let $x$ be a placement position with no captures in which the tigers stand on $T$. Let $\Delta(T)$ be the
smallest total graph distance from the starting squares $\{0, 3, 4\}$ to $T$ over all ways of matching the
three tigers to the three squares. If $\Delta(T)$ exceeds the number of tiger moves made before $x$, then $x$
is unreachable.
\end{lemma}

\begin{proof}
With no capture, every tiger move so far was a slide along one edge, which changes the total distance
travelled by 1. Reaching $T$ needs at least $\Delta(T)$ slides.
\end{proof}

For $T = \{4, 12, 13\}$, $\Delta(T) = 7$. With 4 goats placed and the goats to move, the tigers have made 4
moves if the goats moved first and 5 if the tigers did. The other three positions fail in the same way
($\Delta = 7$ against at most 6 moves).

The tables hold every arrangement of pieces because the recursion needs no more than that, and a position's
value never depends on how it arose. Every stored value is therefore correct. But a maximum taken over all
arrangements is a statement about arrangements, not about games.

\subsection{Reachable positions}

\begin{definition}
For a first mover $f \in \{\gt, \tg\}$, let $x_0^f$ be the opening position (tigers on $\{0, 3, 4\}$, no goats,
$f$ to move). $\Reach_f$ is the smallest set that contains $x_0^f$ and, with every position, all its
successors. For $x \in \Reach_f$, $\delta_f(x)$ is the fewest plies in which a game from $x_0^f$ reaches $x$.
\end{definition}

$\Reach_f$ is computed by a forward search over the same layers and with the same move code as the solver
(\code{reachable.py}; Algorithm~\ref{alg:reach}).

\begin{algorithm}[H]
\caption{Reachable positions and their fewest plies, for first mover $f$}\label{alg:reach}
\begin{tabbing}
\hspace{1.6em}\=\hspace{1.6em}\=\hspace{1.6em}\=\kill
mark $x_0^f$ \\
\textbf{placement, forwards:} for the layers in the order $\Pi^{\tg}_0, \Pi^{\gt}_0, \Pi^{\tg}_1, \Pi^{\gt}_1, \dots, \Pi^{\gt}_{14}$: \\
\> for every marked position of the layer, mark all its successors \\
\textbf{movement, breadth-first:} the successors of $\Pi^{\gt}_{14}$ are the entry positions; set $\delta_f = t_0$ on them \\
\> \textbf{for} $j = t_0, t_0 + 1, \dots$ until nothing new is found: \\
\>\> for every position with $\delta_f = j$ and every successor $y$ without a value: $\delta_f(y) \leftarrow j + 1$
\end{tabbing}
\end{algorithm}

\begin{proposition}\label{prop:reach}
Algorithm~\ref{alg:reach} marks exactly $\Reach_f$, and its $\delta_f$ is the fewest plies from the opening.
In the placement phase the number of plies is the same along every route:
$\delta_f(x) = 2n$ (goats to move) or $2n - 1$ (tigers to move) for $f = \gt$, and $2n + 1$ or $2n$ for $f = \tg$.
\end{proposition}

\begin{proof}
Every placement move goes from a layer to a later one in the stated order (Section~\ref{sec:dp-graph}). When a
layer is processed, all its predecessors' layers are finished, so its marks are final; by induction along
the order the marked set is closed under successors and contains only positions with a marked predecessor
chain to $x_0^f$. The ply count is fixed because each goat turn places exactly one goat and the sides
alternate. In the movement phase every entry position has the same $\delta_f = t_0$ (29 for $f = \gt$, 30 for
$f = \tg$), and breadth-first search from a set of equal depth visits positions in order of distance, so the
first value written is the minimum. The search stops when a level adds nothing, at which point the marked
set is closed under successors.
\end{proof}

Table~\ref{tab:reach} gives the counts. Of the 19.02 billion positions, 18{,}189{,}005{,}282 are reachable when
the goats move first and 18{,}553{,}660{,}020 when the tigers do. Early layers are sparse: after three goats
only 153{,}510 of 2{,}392{,}621 goat-to-move positions can have arisen with the goats first. By the end of
the placement phase, every board with the tigers to move is reachable. The search needs only one more ply to
close: the positions first reached in the movement phase all appear at ply 30 ($f = \gt$) or 31
($f = \tg$). About 1\% of movement positions with the goats to move are never reached; these are boards on
which no tiger can have just moved.

\begin{table}[H]
\centering\scriptsize
\caption{Reachable positions by layer. ``All'' counts every arrangement, per side to move.}\label{tab:reach}
\begin{tabular}{@{}lr rr rr@{}}
\toprule
 & & \multicolumn{2}{c}{goats move first} & \multicolumn{2}{c}{tigers move first} \\
\cmidrule(lr){3-4}\cmidrule(l){5-6}
layer ($n$ placed) & all & goats to move & tigers to move & goats to move & tigers to move \\
\midrule
\tabin{generated/tab_reachable}
\bottomrule
\end{tabular}
\end{table}

\subsection{The longest forced wins of real games}

\begin{definition}
For a first mover $f$ and a winner $w \in \{\gt, \tg\}$,
\begin{align*}
\ell^*_f(w) &= \max\{\, d(x) : x \in \Reach_f,\ w \text{ wins at } x \,\}, \\
\ell^{200}_f(w) &= \max\{\, d(x) : x \in \Reach_f,\ w \text{ wins at } x,\ \delta_f(x) + d(x) \le 200 \,\}.
\end{align*}
\end{definition}

The second quantity is the longest forced win that the 200-ply limit of \code{TnGEnv} does not cut off. By
Proposition~\ref{prop:limit} a win from $x$ survives the limit exactly when the game reached $x$ early enough,
and the earliest it can do so is $\delta_f(x)$.

\begin{theorem}[Longest forced wins]\label{thm:longest}
\begin{center}
\begin{tabular}{@{}lcccc@{}}
\toprule
first mover $f$ & $\ell^*_f(\text{goats})$ & $\ell^{200}_f(\text{goats})$ & $\ell^*_f(\text{tigers})$ & $\ell^{200}_f(\text{tigers})$ \\
\midrule
goats & \textbf{213} & 191 & 67 & 67 \\
tigers & \textbf{214} & 193 & 67 & 67 \\
\bottomrule
\end{tabular}
\end{center}
\end{theorem}

\begin{proof}
The claim has three parts, and each is established by something checked exhaustively.

\emph{(i) The distances are exact.} $d(x)$ in the tables is the true distance for every position, by
Theorem~\ref{thm:certificate} and the full check of Section~\ref{sec:dp-proof}.

\emph{(ii) The set is exact.} The search marks exactly $\Reach_f$ with the true $\delta_f$, by
Proposition~\ref{prop:reach}. Its move generation is the solver's, which equals the environment's on every
board (Section~\ref{sec:dp-proof}).

\emph{(iii) The maximum is taken over all of it.} The program reads $d(x)$ for every marked position and
keeps the largest per winner, and the largest among those with $\delta_f(x) + d(x) \le 200$. Nothing is
sampled. This gives the upper bound: no reachable position has a longer forced win.

\emph{The bound is attained.} For each entry the program outputs a complete game. It walks back from the
record position $x^*$ through marked predecessors to $x_0^f$, which exhibits a legal route of length
$\delta_f(x^*)$, and then plays the perfect line from $x^*$. Each game has been replayed move by move with a
separate implementation of the rules: every move is legal, $x^*$ appears at ply $\delta_f(x^*)$, and the game
ends with the stated winner exactly $d(x^*)$ plies later (Appendix~\ref{app:games}).
\end{proof}

\begin{table}[H]
\centering\small
\caption{The record positions. ``Game'' is the length of the witness game: $\delta_f(x^*) + d(x^*)$.}\label{tab:records}
\begin{tabular}{@{}llp{0.36\linewidth}rrr@{}}
\toprule
$f$ & quantity & position $x^*$ & $\delta_f$ & $d$ & game \\
\midrule
goats & $\ell^*(\text{goats})$ & tigers 1, 5, 8; goats 0, 6, 15, 18, 22; goats to move & 10 & 213 & 223 \\
goats & $\ell^{200}(\text{goats})$ & tigers 4, 5, 8; goats 1, 2, 17; goats to move & 6 & 191 & 197 \\
goats & $\ell^*(\text{tigers})$ & tigers 1, 5, 13; goats 6, 7, 8, 10, 11, 12, 18; 3 captured; tigers to move & 19 & 67 & 86 \\
tigers & $\ell^*(\text{goats})$ & tigers 1, 6, 14; goats 10, 16, 18--22; tigers to move & 14 & 214 & 228 \\
tigers & $\ell^{200}(\text{goats})$ & tigers 0, 5, 7; goats 3, 4, 19; goats to move & 7 & 193 & 200 \\
tigers & $\ell^*(\text{tigers})$ & tigers 1, 5, 13; goats 6, 7, 8, 10, 11, 12, 18; 3 captured; tigers to move & 20 & 67 & 87 \\
\bottomrule
\end{tabular}
\end{table}

\begin{figure}[H]
\centering
\begin{tabular}{ccc}
\tngboard{1,5,8}{0,6,15,18,22}{0.78} & \tngboard{1,6,14}{10,16,18,19,20,21,22}{0.78} & \tngboard{4,12,13}{3,5,6,22}{0.78} \\
\footnotesize goats first: goats to move, & \footnotesize tigers first: tigers to move, & \footnotesize unreachable: goats to move, \\
\footnotesize goats win in 213 & \footnotesize goats win in 214 & \footnotesize goats win in 215
\end{tabular}
\caption{The two record positions of real games, and one of the four positions with the largest distance
over all arrangements. Filled circles are tigers, open circles goats.}\label{fig:records}
\end{figure}

\paragraph{Reading the result.}
\begin{itemize}
\item The reachable records are only one and two plies short of the figure over all arrangements. The
  unreachable positions were not outliers; they sit at the end of a family of very long goat wins.
\item Both records are goat wins that begin early in the placement phase and end with the tigers
  immobilized. The goats' task is slow: they must place eleven or ten more goats and then close the net.
\item The tigers' wins are much shorter. A tiger win is a sequence of captures, and the longest one a game
  can reach takes 67 plies from a position in which three goats are already gone.
\item Both record wins end after ply 200 (at plies 223 and 228 by the shortest route). Under the limit of
  \code{TnGEnv} they would be cut off as draws. The longest win that fits takes 191 plies (goats first) or
  193 (tigers first, ending exactly at ply 200).
\item In the movement phase the largest distance is 210 plies (tigers to move, 15 goats, no captures). That
  position is reachable, but only from ply 29, so its game would last at least 239 plies.
\end{itemize}

\paragraph{Independent checks of this section.} A separate search written in plain Python, sharing no code
with the GPU search, reproduces the reachable counts exactly for the layers with up to four goats placed, for
both first movers (20 layer counts, from 1 to 3{,}295{,}989 positions). All six witness games were replayed
with the visualizer's rule implementation, which is itself tested against \code{TnGEnv}.
